
Filesystems are the organisational structures that determine how data is stored and retrieved on a storage device \parencite{LinuxSAGFilesystems}. They define the way files are named, stored, and accessed, as well as how metadata is managed. Common filesystems in Linux include ext4, XFS, Btrfs, and F2FS, each with its own features and performance characteristics.

The Linux filesystem provides the fundamental infrastructure through which processes interact with persistent data and communicate with each other. Understanding file operations is crucial for analysing how AI agents might manipulate system state, hide information, or detect the presence of other agents.

\subsection{Data Blocks, Clusters and Sectors}
\label{DataBlocks}

Sectors refer to the smallest physical unit of storage on a traditional disc drive, typically 512 bytes of storage. Sectors used to be fixed in size due to the limitations and design of hard disc drives, but with modern solid-state drives the notion of sectors has become a logical abstraction rather than a physical constraint.

When a collection of sectors is grouped together, it forms a block or cluster depending on the filesystem. The concept of a block as a machine data unit was formalised in the design of the IBM 7030 (Stretch) in the early 1960s, where it denoted the number of words transmitted to or from an input-output unit in response to a single I/O instruction, distinguishing the machine data hierarchy (bit, byte, word, block) from the natural data hierarchy of the data being processed (character, field, record, file) \parencite{ProjectStretch}. A block is a logical unit of storage that the filesystem uses to manage files. In ext4, a block is a group of sectors whose size must be an integral power of two, ranging from 1\,KiB to 64\,KiB; the block size is fixed at \texttt{mkfs} time and is typically 4\,KiB \parencite{KernelExt4Blocks}. Blocks are in turn grouped into larger units called block groups, which allow the filesystem to localise metadata and data for related files \parencite{KernelExt4Blocks}. Clusters are essentially the same concept as blocks, but the term is more commonly used in the context of NTFS (New Technology File System), the default filesystem for modern Windows operating systems \parencite{MicrosoftNTFS}.

Blocks are the smallest logical unit of storage, meaning that even if a file is only a few bytes in size, it will still occupy at least one block on the filesystem. This can lead to inefficiencies, especially for small files, as the unused space within the block cannot be utilised by other files. The choice of block size affects both the maximum addressable filesystem size and the maximum file size: at 4\,KiB blocks, a 32-bit ext4 filesystem can address up to 16\,TiB, while enabling the 64-bit feature extends this to 64\,ZiB \parencite{KernelExt4Blocks}.

\subsubsection{Super Blocks}
\label{SuperBlocks}

Blocks are tracked by the filesystem's metadata structures, such as inodes or the superblock, which record which blocks are allocated to which files. In the VFS object model, a superblock object represents a mounted filesystem instance \parencite{KernelVFS}. In ext4, the superblock is a 1024-byte on-disc structure that records global filesystem information: total and free block counts, total and free inode counts, the block size, the filesystem UUID, the volume label, mount timestamps, error statistics, and the set of compatible, incompatible, and read-only-compatible feature flags that govern how the kernel and \texttt{fsck} may interact with the filesystem \parencite{KernelExt4Super}. If the \texttt{sparse\_super} feature flag is set, redundant copies of the superblock are kept only in block groups whose number is 0 or a power of 3, 5, or 7, rather than in every group, reducing metadata overhead \parencite{KernelExt4Super}. When a file is created or modified, the filesystem allocates blocks to store the file's data and updates the metadata accordingly. The block size can impact performance and storage efficiency; larger blocks can reduce fragmentation but may waste space for small files, while smaller blocks can be more efficient for small files but may lead to increased fragmentation.

\subsubsection{Block Groups}
\label{BlockGroups}

Blocks are organised into block groups, each of which bundles together the metadata and data for a contiguous region of the filesystem. The layout of a standard block group begins with a copy of the superblock and group descriptor table (present only in selected groups), followed by reserved GDT blocks for future expansion, a data block bitmap, an inode bitmap, the inode table, and finally the actual data blocks \parencite{KernelExt4BlockGroup}. Block group 0 is a special case: the first 1024 bytes are left unused to accommodate x86 boot sectors, after which the superblock begins \parencite{KernelExt4BlockGroup}.

The primary superblock and group descriptor table reside in block group 0; redundant copies are spread across other groups so that corruption at the start of the disc does not render the entire filesystem unreadable. Reserved GDT block space is also pre-allocated to allow the filesystem to grow by up to 1024$\times$ its original size without reformatting \parencite{KernelExt4BlockGroup}.

Ext4 introduces flexible block groups (\texttt{flex\_bg}), in which several consecutive block groups are treated as a single logical unit. The bitmaps and inode tables of all constituent groups are consolidated into the first group of the flex\_bg, keeping metadata contiguous on disc and enabling large files to be stored in a continuous region \parencite{KernelExt4BlockGroup}. For very large filesystems, the meta block group feature (META\_BG) relocates group descriptor copies from the congested first block group into the first, second, and last groups of each metablock group, raising the maximum number of block groups from $2^{21}$ to the hard limit of $2^{32}$ and supporting filesystems up to 512\,PiB \parencite{KernelExt4BlockGroup}.

\subsubsection{Free Space Management}
\label{FreeSpaceManagement}

The filesystem must track which blocks are free and which are allocated in order to service file creation and extension requests. Four principal strategies exist for maintaining this free space information \parencite{GFGFreeSpaceManagement,GFGOSFileSystems}. The bitmap (or bit vector) approach assigns one bit to each block: a zero bit indicates a free block and a one bit indicates an allocated block; finding the first free block reduces to scanning for the first non-zero word in the bitmap, after which the first set bit within that word identifies the block. This is the strategy used by ext4, whose data block bitmap and inode bitmap are stored at the start of each block group (see \autoref{BlockGroups}). The linked list approach chains free blocks together so that each free block stores the address of the next; this avoids dedicated bitmap storage but requires disc I/O to traverse the list. The boundary tag approach embeds a size and allocation status marker in each block header, which simplifies coalescing adjacent free blocks during deallocation and reduces fragmentation. Finally, a free list approach maintains an explicit array or linked list of pointers to free blocks directly in memory, enabling fast allocation at the cost of additional memory for the list itself \parencite{GFGFreeSpaceManagement}.

\subsubsection{Boot Blocks}
\label{BootBlocks}

The boot block (or boot sector) is a reserved area at the beginning of a storage device or partition that contains the code necessary to start the boot process. On traditional BIOS systems with MBR (Master Boot Record) partitioning, the boot block occupies the first 512 bytes of the disc and contains both the bootloader code and the partition table. The bootloader is a small program that the system firmware executes to locate and load the operating system kernel into memory.

On modern UEFI (Unified Extensible Firmware Interface) systems with GPT (GUID Partition Table), the boot process is more complex and uses an EFI System Partition (ESP), which is a dedicated FAT32 partition containing bootloader files. However, GPT discs still maintain a protective MBR in the first sector for backward compatibility with legacy systems.

Boot blocks are critical for system initialisation and are protected areas that typically cannot be modified by normal filesystem operations. In multi-boot scenarios, boot managers like GRUB (Grand Unified Bootloader) or systemd-boot reside in these boot blocks and provide menus for selecting which operating system to load. Understanding boot blocks is important in security contexts, as malicious actors might attempt to modify boot code to achieve persistence or execute code before the operating system loads (bootkits or rootkits).

\subsection{Inodes}
\label{Inodes}

At the filesystem implementation level, files are represented by inodes (index nodes), which are data structures used by many Unix-like filesystems to represent files and directories \parencite{KernelVFS}. Each inode contains metadata about a file, including its size, permissions, ownership, timestamps (creation, modification, access), and pointers to the actual data blocks where the file's content is stored. The inode does not contain the filename itself; instead, directory entries map filenames to their corresponding inode numbers. The inode contains the numbers of the data blocks where the file's content is stored; when a file is large enough that more pointers are needed than can fit directly in the inode, additional indirect blocks are dynamically allocated to hold the overflow pointers \parencite{LinuxSAGFilesystems}. Inodes for block device filesystems reside on disc and are copied into memory when required; inodes for pseudo-filesystems such as \texttt{/proc} exist only in memory \parencite{KernelVFS,KernelProc}.

When a file is accessed, the filesystem uses the directory structure to resolve the filename to an inode number, then retrieves the inode to access the file's metadata and data blocks. This separation of filename and file metadata allows for features like hard links, where multiple directory entries can point to the same inode, effectively giving a single file multiple names.

\subsection{Filesystems and The Virtual File System}
\label{VFS}

Linux organises storage into three conceptual levels: the virtual file system (VFS), individual filesystems, and files themselves \parencite{KernelVFS,GFGOSFileSystems}. The Virtual File System (also known as the Virtual Filesystem Switch) is a software layer within the Linux kernel that provides the filesystem interface to userspace programs and an abstraction that allows different filesystem implementations to coexist \parencite{KernelVFS}. It allows the kernel to support multiple filesystem types without requiring changes to userspace applications, and ensures that standard system calls such as \texttt{open()}, \texttt{read()}, \texttt{write()}, and \texttt{stat()} behave consistently regardless of the underlying filesystem type.

When a pathname is passed to one of these system calls, the VFS translates it using its directory entry cache (also known as the dentry cache, or dcache). The dcache provides a fast look-up mechanism to translate pathnames into dentries, which in turn point to the inodes that hold the file's metadata and data block pointers \parencite{KernelVFS}.

At the top level, Linux employs a single hierarchical directory tree\parencite{GFGUnixFileSystem} rooted at \texttt{/}, unlike Windows which uses drive letters. All storage devices, regardless of their physical location or filesystem type, are mounted into this unified tree structure. Common mount points include \texttt{/home} for user directories, \texttt{/tmp} for temporary files, and \texttt{/proc} or \texttt{/sys} for kernel-provided pseudo-filesystems that expose process and system information \parencite{KernelProc}. The standard Unix directory hierarchy also includes \texttt{/bin} and \texttt{/usr/bin} for executable programs, \texttt{/etc} for system-wide configuration files, \texttt{/dev} for device file representations, \texttt{/lib} for system libraries, \texttt{/var} for variable data such as logs and mail spools, and \texttt{/root} as the home directory for the system administrator \parencite{GFGUnixFileSystem}.

\subsubsection{Filesystem Types}
\label{FilesystemTypes}

Individual filesystems can be of various types, each with distinct characteristics. Modern Linux distributions commonly use ext4, which provides journaling for crash recovery and supports large files and volumes. Journaling means that a record is kept of filesystem transactions before they are committed, allowing the filesystem to reconstruct a consistent state after an unclean shutdown without requiring a full disc check \parencite{LinuxSAGFilesystems}. The ext (extended) filesystem family, including ext2, ext3, and ext4, is one of the most widely used filesystem families in Linux \parencite{TLDPFilesystems}. Ext2 introduced the block group layout still used by ext4 today, and provides standard Unix file semantics including long filenames of up to 255 characters. By default, ext2 reserves 5\% of disc blocks for the superuser, allowing an administrator to recover from situations where user processes have filled the filesystem entirely \parencite{TLDPFilesystems}. Short symbolic links whose target path fits within 60 characters are stored directly inside the inode rather than in a separate data block, saving disc space and eliminating an extra I/O operation on access \parencite{TLDPFilesystems}. Ext2 also tracks filesystem state in the superblock: the kernel marks the filesystem as ``Not Clean'' when it is mounted read/write and resets it to ``Clean'' on unmount, so that \texttt{fsck} can skip the consistency check when the state is clean; if the kernel detects an inconsistency it marks the filesystem as ``Erroneous'' to force a check at the next boot \parencite{TLDPFilesystems}. Ext3 extended ext2 with journaling for crash recovery. Ext4 is the current default, adding extents-based allocation, 64-bit addressing, and further performance improvements.

Alternative filesystems include XFS (optimised for parallel I/O operations), Btrfs (supporting snapshots and copy-on-write), and F2FS (optimised for flash storage). In the Windows ecosystem, NTFS is the default filesystem, offering transaction-based journaling, granular access control via ACLs, and support for very large volumes (up to 8\,PB with a 2\,MiB cluster size) \parencite{MicrosoftNTFS}. The FAT (File Allocation Table) filesystem family, although older, remains widely used on removable media such as USB drives due to its broad compatibility across operating systems. On macOS and iOS devices, Apple's HFS+ (Hierarchical File System Plus) and its successor APFS (Apple File System) are used, with APFS introducing features such as copy-on-write cloning, snapshots, and native encryption \parencite{GFGOSFileSystems}. The choice of filesystem affects performance characteristics, maximum file sizes, and available features, though the VFS layer ensures that basic operations remain consistent across types \parencite{KernelFilesystems}.

\subsection{Partitions}
\label{Partitions}

Partitions are divisions of a storage device that allow it to be managed separately. Each partition can contain a different filesystem, enabling multiple operating systems or different types of data to coexist on the same physical device. Partitions are defined in a partition table, which is typically located at the beginning of the storage device. Common partitioning schemes include MBR (Master Boot Record) and GPT (GUID Partition Table).

\subsection{File Operations}
\label{FileOperations}

Files can be created, modified, copied, moved, and deleted using a variety of commands. The \texttt{touch} command creates empty files or updates timestamps, while \texttt{mkdir} creates directories. File creation operations (including creating a new file via \texttt{open()} with \texttt{O\_CREAT} and creating a directory with \texttt{mkdir}) are serialised by the kernel through exclusive acquisition of the parent directory's \texttt{i\_rwsem} \parencite{KernelFSLocking}. File content can be written using redirection operators (\texttt{>} and \texttt{>>}) or commands like \texttt{echo}, \texttt{cat}, or text editors.

Copying and moving files is accomplished through \texttt{cp} and \texttt{mv} respectively. The \texttt{cp} command creates a new file with duplicated content, while \texttt{mv} relocates the file within the filesystem hierarchy without duplicating data when the source and destination are on the same filesystem. Deletion is performed with \texttt{rm} for files and \texttt{rmdir} or \texttt{rm -r} for directories.

More sophisticated operations include \texttt{rsync}, which can synchronise files between directories or systems while minimising data transfer through checksums and delta encoding, and \texttt{dd}, which performs low-level block-by-block copying useful for creating disc images or manipulating raw devices.

\subsection{Links}
\label{Links}

This separation of filename and inode enables the concept of links. A hard link is an additional directory entry pointing to the same inode, effectively giving a single file multiple names. All hard links are equivalent, and the file's data is only deleted when the last link is removed and no processes have the file open \parencite{KernelVFS}. Hard links cannot span filesystems and cannot reference directories (to prevent circular references). At the kernel level, the \texttt{link()} and \texttt{unlink()} inode operations both acquire an exclusive lock on the affected inodes' read/write semaphore (\texttt{i\_rwsem}), serialising concurrent modifications to directory entries and inode link counts \parencite{KernelFSLocking}.

Symbolic links (symlinks), created with \texttt{ln -s}, are special files containing a path to another file. Unlike hard links, symlinks can reference directories and cross filesystem boundaries, but they break if the target is moved or deleted. The \texttt{readlink} command can resolve symbolic links to their targets.

Beyond regular files and directories, Unix filesystems recognise several additional file types, each represented by a distinct entry in the directory hierarchy \parencite{GFGUnixFileSystem}. Character special files and block special files, located under \texttt{/dev}, represent physical and virtual devices: character devices transfer data one byte at a time (such as keyboard or serial port devices), while block devices transfer data in fixed-size blocks (such as hard discs and SSDs). Named pipes (FIFOs) provide a one-way inter-process data channel without a network connection. Unix domain sockets support bidirectional inter-process communication entirely within the local filesystem. All types are reported by \texttt{ls -l} using a character in the first column of the output: \texttt{-} for regular files, \texttt{d} for directories, \texttt{c} for character special files, \texttt{b} for block special files, \texttt{p} for named pipes, \texttt{s} for sockets, and \texttt{l} for symbolic links \parencite{GFGUnixFileSystem}.

\subsection{File Descriptors}
\label{FileDescriptors}

File descriptors are integer handles that processes use to reference open files. When a process opens a file using the \texttt{open()} system call, the VFS allocates a file structure, which is the kernel-side implementation of a file descriptor \parencite{KernelVFS}. This file structure is initialised with a pointer to the relevant dentry and a set of file operation functions taken from the inode, and is then placed into the file descriptor table for the process. The kernel returns the integer index into that table as the file descriptor. Standard file descriptors include 0 (stdin), 1 (stdout), and 2 (stderr). Commands like \texttt{lsof} (list open files) can display which file descriptors are in use by which processes, making it possible to detect which files a user might be accessing.

\subsection{Permissions and Access Control}
\label{PermissionsAndACL}

Linux implements a permission model that controls read, write, and execute access for three categories: the file's owner, members of the file's group, and all other users. Permissions are displayed by \texttt{ls -l} in the format \texttt{rwxrwxrwx}, where each trio represents owner, group, and other permissions respectively.

The \texttt{chmod} command modifies permissions using either symbolic notation (\texttt{chmod u+x file}) or octal notation (\texttt{chmod 755 file}). At the kernel level, the \texttt{setattr} inode operation that underlies \texttt{chmod} and \texttt{chown} acquires an exclusive \texttt{i\_rwsem} lock on the inode, serialising permission changes \parencite{KernelFSLocking}. Conversely, permission checks may be performed in RCU-walk mode without holding any lock, allowing the kernel to validate access rights on the fast path \parencite{KernelFSLocking}. Ownership is changed with \texttt{chown} for the owner and \texttt{chgrp} for the group. The special permissions bits include the setuid bit (executes with owner's privileges), setgid bit (inherits directory's group), and sticky bit (restricts deletion in shared directories like \texttt{/tmp}).

Beyond traditional Unix permissions, modern Linux systems support Access Control Lists (ACLs), which provide finer-grained control by allowing permissions to be specified for arbitrary users and groups. ACLs are managed through \texttt{getfacl} and \texttt{setfacl} commands. Additionally, Security-Enhanced Linux (SELinux) and AppArmor provide mandatory access control frameworks that can restrict processes regardless of traditional file permissions.

File attributes, accessible through \texttt{lsattr} and modifiable via \texttt{chattr}, provide additional control mechanisms. For example, the immutable attribute (\texttt{+i}) prevents any modifications to a file, even by root, until the attribute is removed.

\subsection{Mounting and Filesystem Management}
\label{MountingFilesystems}

Filesystems must be mounted before they can be accessed, using the \texttt{mount} command which attaches a filesystem to a directory mount point. In the kernel, mounting is performed through a multi-step process: a filesystem context (\texttt{fs\_context}) is created, mount parameters are parsed and attached to it, the context is validated, a superblock is obtained or created, and the mount is performed before the context is destroyed \parencite{KernelMountAPI}. The \texttt{fs\_context} structure holds the mounter's credentials, the source device or path, namespace references, security data, and the superblock flags that will govern the mounted filesystem \parencite{KernelMountAPI}. This separation of parameter collection from superblock creation means that individual mount options can be passed and validated independently before the filesystem is committed, rather than being handed to the kernel as a single opaque data page as in the legacy \texttt{mount(2)} interface \parencite{KernelMountAPI}. The \texttt{umount} command detaches filesystems. The \texttt{/etc/fstab} file specifies which filesystems should be automatically mounted at boot time.

Creating new filesystems is accomplished with the \texttt{mkfs} family of commands (e.g., \texttt{mkfs.ext4}, \texttt{mkfs.xfs}), which initialise the bookkeeping data structures on a partition before it can be used \parencite{LinuxSAGFilesystems}. Filesystem consistency checks and repairs are performed by \texttt{fsck}, which must only be run on unmounted filesystems, since accessing the raw disc while the operating system also has the filesystem mounted will cause inconsistencies \parencite{LinuxSAGFilesystems}. These operations typically require elevated privileges.

\subsection{Monitoring Filesystem Usage and Activity}
\label{MonitoringFilesystem}

Monitoring filesystem usage is essential for detecting unusual behaviour. The \texttt{df} command displays disc space usage by filesystem, while \texttt{du} shows space used by directories and files; both are useful for locating disk space consumers \parencite{LinuxSAGFilesystems}. The \texttt{stat} command provides detailed information about a specific file, including inode number, size, blocks allocated, and all timestamps (access time, modification time, and change time).

The \texttt{/proc} pseudo-filesystem is a particularly rich source of live system and per-process information \parencite{KernelProc}. Each running process has a subdirectory at \texttt{/proc/\textit{PID}/} containing files such as \texttt{status} (human-readable process state, memory usage, and signal masks), \texttt{fd/} (symbolic links to all open file descriptors), and \texttt{maps} (memory region mappings with access permissions). Utilities such as \texttt{ps} obtain their process information from \texttt{/proc} \parencite{KernelProc}. Accessing another process's \texttt{/proc/\textit{PID}/} entries requires either the \texttt{CAP\_SYS\_PTRACE} capability with \texttt{PTRACE\_MODE\_READ} permissions, or the \texttt{CAP\_PERFMON} capability \parencite{KernelProc}.

For detecting file system activity, several commands prove useful. The \texttt{find} command can search for files based on numerous criteria including name patterns, modification times, permissions, and size. The \texttt{inotify} system provides efficient filesystem event monitoring, accessible through tools like \texttt{inotifywait}, which can watch for file creation, modification, deletion, and access in real time.

